% Claim statement file for row claim_3_5 -- "BifurcationAlternative".
% Auto-generated stub by scaffold/solve_chapter.py at row start. The
% manager agent (or a worker it dispatches) fills the body below with the
% claim's statement(s) from the lecture notes -- statement only, no proof
% (proofs live in the sibling `claim_3_5_proof_*.tex` file).
%
% This file is a sibling of `main.tex` inside `<subsection>/tex/`. The
% `subfiles` package lets it render both standalone and as part of
% `main.tex`. Note: `main.tex` does NOT \subfile this statement file -- the
% sibling `_proof_` file restates the statement above the proof, so
% including both in `main.tex` would be redundant. This file exists for
% standalone reading / grep convenience.

\documentclass[main]{subfiles}

\begin{document}
% ref: claim_3_5 (statement)
% title: BifurcationAlternative
\def\rowref{claim\_3\_5}
\phantomsection\label{claim_3_5}

% NOTE on the bifurcation-source convention assumed by this proposition.
%
% Literal LN def 3.4 (Walks, item vi.) writes a bifurcation as
%     $v = v_0 \hut v_1 \hut \cdots \hut v_{k-1} \hus v_k \tuh \cdots \tuh v_n = w$
% with $1 \le k \le n$ and source $v_k$ at the directed hinge. The boundary
% case $k = n$ with directed hinge (right arm trivial, single edge
% $w \to v_{n-1}$) would let the source coincide with $w$, contradicting the
% RHS clause $c \in \Anc^{G_{\doit(v)}}(w) \sm \{w\}$ in any acyclic graph
% (and in particular whenever $w \in J$, since by claim \ref{claim_3_1}
% input nodes have no incoming edges, so $\Anc^{G_{\doit(v)}}(w) = \{w\}$).
%
% Our def_3_4 encoding has already committed to the strictly-interior-source
% convention via the chapter-init addition recorded in clause (e) of
% def \ref{def:walks}, item~vi. ("each end-node has exactly one arrowhead
% pointing toward it"), which excludes the $k = n$ directed-hinge case. As
% stated in the closing paragraph "Source of the bifurcation" of
% def \ref{def:walks}, item~vi., whenever a source $v_k$ is defined we
% automatically have $1 \le k \le n - 1$, so the source is an interior
% vertex distinct from both end-nodes $v_0 = v$ and $v_n = w$. Under this
% convention the LN proposition below is provable as-stated, and no
% addition_to_the_LN clause is needed for claim_3_5.

\begin{Prp}[BifurcationAlternative]\label{prp:bifurcations_alternative}
    Let $G = (J, V, E, L)$ be a CDMG (def \ref{def-cdmg}); throughout we
    use the notation of def \ref{not-cdmg}, the walk definitions of def
    \ref{def:walks}, the family-relationship notation of def
    \ref{def_3_5}, and the hard-intervention construction of def
    \ref{def:G_hard_intervention}. Let
    \[ v,\ w,\ c \;\in\; J \cup V \]
    be three (not necessarily distinct) nodes. All four quantifiers
    ``for every CDMG $G$'', ``for every $v \in J \cup V$'', ``for every
    $w \in J \cup V$'', and ``for every $c \in J \cup V$'' are read
    \emph{universally}, and the biconditional below is asserted for every
    such tuple $(G, v, w, c)$.

    \emph{Hard interventions on singletons.}
    The expressions $G_{\doit(w)}$ and $G_{\doit(v)}$ appearing on the
    right-hand side are read as the singleton-set hard interventions
    $G_{\doit(\{w\})}$ and $G_{\doit(\{v\})}$ in the sense of def
    \ref{def:G_hard_intervention} (the LN's lowercase $w$, $v$ inside the
    $\doit(\cdot)$ slot is shorthand for the singleton sets $\{w\}$ and
    $\{v\}$). Since $\{w\}, \{v\} \subseteq J \cup V$ the precondition of
    def \ref{def:G_hard_intervention} is met, so both intervened CDMGs
    are defined; and by items i.\ and ii.\ of def
    \ref{def:G_hard_intervention},
    \[
        J_{\doit(\{w\})} \cup V_{\doit(\{w\})}
        \;=\; (J \cup \{w\}) \cup (V \sm \{w\}) \;=\; J \cup V,
    \]
    and analogously for $\{v\}$ in place of $\{w\}$. Consequently
    $\Anc^{G_{\doit(w)}}(\cdot)$ and $\Anc^{G_{\doit(v)}}(\cdot)$ (def
    \ref{def_3_5}, item~iv.) are evaluated on the same carrier $J \cup V$
    on which $v$, $w$, $c$ range, and the set-minus clauses
    $\sm \{v\}$ and $\sm \{w\}$ are well-typed.

    \emph{The biconditional.} Propositions (a) and (b) below are
    equivalent.
    \begin{enumerate}[label=(\alph*)]
        \item \emph{Existence of a bifurcation between $v$ and $w$ with
              source $c$.}
              There exists a walk
              \[
                  p \;=\; (v_0, a_0, v_1, a_1, \dots, v_{n-1}, a_{n-1}, v_n)
              \]
              in $G$ from $v$ to $w$ (so $v_0 = v$ and $v_n = w$, in the
              sense of def \ref{def:walks}, item~i., with some integer
              $n \ge 0$), together with an index $k$ with
              $1 \le k \le n$, such that both of the following hold:
              \begin{itemize}
                  \item $p$ is a bifurcation between $v$ and $w$ in $G$
                        with witness index $k$, in the sense of def
                        \ref{def:walks}, item~vi., i.e.\ all five
                        clauses (a)--(e) of def \ref{def:walks},
                        item~vi.\ hold for $p$ at the index $k$;
                  \item $c$ is a source of this bifurcation, in the sense
                        of def \ref{def:walks}, item~vi., closing
                        paragraph ``Source of the bifurcation'', i.e.\
                        clause~(d) of def \ref{def:walks}, item~vi.\ is
                        realised at $k$ by the \emph{directed}
                        alternative $a_{k-1} = (v_k, v_{k-1}) \in E$,
                        and $c = v_k$.
              \end{itemize}
              By clause~(e) of def \ref{def:walks}, item~vi., the
              combination of these two requirements forces the witness
              index to satisfy $1 \le k \le n - 1$ (and hence $n \ge 2$);
              so the source $c = v_k$ is automatically an interior vertex
              of $p$, with $c \ne v$ and $c \ne w$.

        \item \emph{Set-theoretic ancestral characterisation.}
              All three of the following clauses hold:
              \begin{enumerate}[label=(\roman*)]
                  \item $v \ne w$;
                  \item $c \in \Anc^{G_{\doit(w)}}(v) \sm \{v\}$;
                        equivalently, unfolding $\Anc$ via def
                        \ref{def_3_5}, item~iv., and the singleton-set
                        reading of $\doit(w)$ above: \emph{$c$ is an
                        ancestor of $v$ in the do-on-$\{w\}$ intervened
                        CDMG $G_{\doit(\{w\})}$}, i.e.\ there exists a
                        directed walk from $c$ to $v$ in
                        $G_{\doit(\{w\})}$ (def \ref{def:walks},
                        item~ii.; def \ref{def:G_hard_intervention}),
                        \emph{and} $c \ne v$;
                  \item $c \in \Anc^{G_{\doit(v)}}(w) \sm \{w\}$;
                        equivalently, unfolding $\Anc$ via def
                        \ref{def_3_5}, item~iv., and the singleton-set
                        reading of $\doit(v)$ above: \emph{$c$ is an
                        ancestor of $w$ in the do-on-$\{v\}$ intervened
                        CDMG $G_{\doit(\{v\})}$}, i.e.\ there exists a
                        directed walk from $c$ to $w$ in
                        $G_{\doit(\{v\})}$ (def \ref{def:walks},
                        item~ii.; def \ref{def:G_hard_intervention}),
                        \emph{and} $c \ne w$.
              \end{enumerate}
              In LN macro form, (b) is the conjunction
              \[
                  v \ne w \;\land\;
                  c \in \Anc^{G_{\doit(w)}}(v) \sm \{v\} \;\land\;
                  c \in \Anc^{G_{\doit(v)}}(w) \sm \{w\},
              \]
              matching the LN's RHS verbatim.
    \end{enumerate}
\end{Prp}

\end{document}
